\documentclass[12pt,a4paper]{article}

%% ════════════════════════════════════════════════════════════════════════════
%%  PAQUETES
%% ════════════════════════════════════════════════════════════════════════════
\usepackage[utf8]{inputenc}
\usepackage[T1]{fontenc}
\usepackage{amsmath,amssymb,amsfonts,amsthm,mathtools,bm}
\usepackage{geometry}\geometry{left=2.5cm,right=2.5cm,top=3cm,bottom=3cm}
\usepackage{booktabs,longtable,multirow,array,tabularx}
\usepackage{graphicx,float,caption}
\usepackage[table,xcdraw]{xcolor}
\usepackage{hyperref}
\hypersetup{colorlinks=true,linkcolor=uptv,urlcolor=uptv,citecolor=uptv}
\usepackage{enumitem,fancyhdr,titlesec,setspace,parskip}
\usepackage{tcolorbox}
\tcbuselibrary{breakable,skins}
\usepackage{listings}
\usepackage{pgfplots,tikz}
\pgfplotsset{compat=1.18}
\usepackage{mdframed}

%% ════════════════════════════════════════════════════════════════════════════
%%  COLORES INSTITUCIONALES UPTC
%% ════════════════════════════════════════════════════════════════════════════
\definecolor{uptv}{RGB}{0,51,102}
\definecolor{oro}{RGB}{184,152,53}
\definecolor{gris_f}{RGB}{247,247,252}
\definecolor{rojo_a}{RGB}{165,20,20}
\definecolor{verde_o}{RGB}{10,105,30}
\definecolor{azul_2}{RGB}{0,80,160}
\definecolor{naranja}{RGB}{200,90,0}
\definecolor{gris_t}{RGB}{90,90,90}
\definecolor{uptc_verde}{RGB}{0,102,51}

%% ════════════════════════════════════════════════════════════════════════════
%%  CABECERA / PIE DE PÁGINA
%% ════════════════════════════════════════════════════════════════════════════
\pagestyle{fancy}\fancyhf{}
\fancyhead[L]{\scriptsize\color{uptv}\textbf{Econometría · UPTC}}
\fancyhead[C]{\scriptsize\color{uptv}Modelo Log-Log · Inflación Colombia 1975--2005}
\fancyhead[R]{\scriptsize\color{uptv}E. Quintana Silva}
\fancyfoot[C]{\small\color{gris_t}\thepage}
\renewcommand\headrulewidth{1pt}
\renewcommand\headrule{\hbox to\headwidth{%
  \color{uptv}\leaders\hrule height\headrulewidth\hfill}}

%% ════════════════════════════════════════════════════════════════════════════
%%  TITULOS DE SECCIÓN
%% ════════════════════════════════════════════════════════════════════════════
\titleformat\section
  {\large\bfseries\color{uptv}}{\thesection.}{.5em}{}
  [\color{uptv}\titlerule]
\titleformat\subsection
  {\normalsize\bfseries\color{uptv}}{\thesubsection}{.5em}{}
\titleformat\subsubsection
  {\small\bfseries\color{uptv!85}}{\thesubsubsection}{.5em}{}

%% ════════════════════════════════════════════════════════════════════════════
%%  CAJAS TCOLORBOX
%% ════════════════════════════════════════════════════════════════════════════
\newtcolorbox{defbox}[1]{
  colback=uptv!6,colframe=uptv,
  fonttitle=\small\bfseries,title={#1},
  breakable,left=4pt,right=4pt,top=2pt,bottom=2pt}

\newtcolorbox{resbox}{
  colback=verde_o!5,colframe=verde_o,
  fonttitle=\small\bfseries,
  title={$\blacktriangleright$\ Resultado (verificado con Stata 17)},
  breakable}

\newtcolorbox{alertbox}{
  colback=rojo_a!5,colframe=rojo_a,
  fonttitle=\small\bfseries,
  title={$\blacktriangle$\ Advertencia Econométrica},
  breakable}

\newtcolorbox{notabox}{
  colback=azul_2!5,
  colframe=azul_2,
  fonttitle=\small\bfseries,
  title={Nota Técnica / Equivalencia Stata--R},
  width=1.2\textwidth, % Aumenta el ancho al 120%
  center, % Centra el cuadro si excede el ancho de texto
  breakable
}

\newtcolorbox{frmbox}[1]{
  colback=oro!8,colframe=oro,
  fonttitle=\small\bfseries,title={#1},
  breakable}

\newtcolorbox{statabox}{
  colback=gray!8,
  colframe=gray!50,
  fonttitle=\small\bfseries\ttfamily,
  title={Salida Stata},
  breakable,
  left=4pt,
  width=1.1\textwidth, % Aumenta el ancho al 110%
  center % Centra el cuadro si excede el ancho del texto
}
%% ════════════════════════════════════════════════════════════════════════════
%%  CÓDIGO
%% ════════════════════════════════════════════════════════════════════════════
\lstdefinestyle{Rcode}{
  language=R,basicstyle=\ttfamily\scriptsize,
  keywordstyle=\color{blue}\bfseries,
  commentstyle=\color{gris_t}\itshape,
  stringstyle=\color{rojo_a!80},
  backgroundcolor=\color{gray!7},
  frame=single,rulecolor=\color{gray!40},
  breaklines=true,numbers=left,
  numberstyle=\tiny\color{gray},
  numbersep=5pt,xleftmargin=15pt,captionpos=b}

\lstdefinestyle{Stata}{
  basicstyle=\ttfamily\small,
  backgroundcolor=\color{gray!7},
  frame=single,rulecolor=\color{gray!40},
  breaklines=true}

%% ════════════════════════════════════════════════════════════════════════════
%%  MACROS MATEMÁTICAS
%% ════════════════════════════════════════════════════════════════════════════
\newcommand\bhat{\hat{\bm\beta}}
\newcommand\X{\mathbf{X}}
\newcommand\y{\mathbf{y}}
\newcommand\Xt{\mathbf{X}'}
\newcommand\E{\mathbb{E}}
\newcommand\Var{\operatorname{Var}}
\newcommand\Cov{\operatorname{Cov}}
\newcommand\YES{\textcolor{verde_o}{\textbf{Sí\,\checkmark}}}
\newcommand\NOO{\textcolor{rojo_a}{\textbf{No\,\texttimes}}}

\setlength\parindent{0pt}
\setlength\parskip{5pt}
\onehalfspacing

%% ════════════════════════════════════════════════════════════════════════════
\begin{document}

%% ─────────────────────────────────────────────────────────────────────────
%%  PORTADA
%% ─────────────────────────────────────────────────────────────────────────
\begin{titlepage}\centering
{\color{uptv}\rule{\linewidth}{3pt}}\vspace{2pt}
{\color{oro}\rule{\linewidth}{1.5pt}}
\vspace{1.6cm}

{\fontsize{24}{30}\selectfont\bfseries\color{uptv}
ANÁLISIS ECONOMÉTRICO\\[0.35em]
MODELO LOG-LOG DE INFLACIÓN}

\vspace{0.5cm}
{\large\color{uptv}\itshape Colombia, período 1975--2005}

\vspace{0.6cm}
{\color{oro}\rule{0.6\linewidth}{0.8pt}}
\vspace{0.6cm}

{\large
Diagnóstico de Multicolinealidad, Sesgo de Especificación\\[0.25em]
y Coherencia Teórica}\\[0.3em]
{\normalsize\color{gris_t}
Resultados verificados con Stata 17 ·
$n=31$ observaciones anuales}

\vspace{2cm}
\begin{tabular}{@{\hspace{8mm}}ll@{}}
\toprule
\rule{0pt}{10pt}\textbf{Estudiante}    & Emanuel Quintana Silva\\[3pt]
\textbf{Institución}   & Universidad Pedagógica y Tecnológica de Colombia (UPTC)\\[3pt]
\textbf{Docente}       & Profesora Rosalba Gil Galindo\\[3pt]
\textbf{Curso}         & Econometría\\[3pt]
\textbf{Modelo}        & Log-Log · MCO · Regresión múltiple\\[3pt]
\textbf{Variable dep.} & Inflación (\%)\\[3pt]
\textbf{Regresores}    & PIB $\cdot$ Tasa de Interés (TI) $\cdot$ Oferta Monetaria (M)\\[3pt]
\textbf{Software}      & Stata 17 + R 4.x (car, lmtest, sandwich)\\[3pt]
\textbf{Fecha}         & \today\\
\bottomrule
\end{tabular}

\vspace{0.8cm}
\begin{mdframed}[linecolor=oro,linewidth=1pt,backgroundcolor=oro!5]
\centering\small\itshape
``Con profundo agradecimiento a la Profesora Rosalba Gil Galindo\\
por su dedicada orientación en el aprendizaje de la Econometría.''
\end{mdframed}

\vspace{1.5cm}
{\color{oro}\rule{\linewidth}{1.5pt}}\vspace{2pt}
{\color{uptv}\rule{\linewidth}{3pt}}
\end{titlepage}

\tableofcontents
\newpage

%% ════════════════════════════════════════════════════════════════════════════
\section{Introducción y contexto macroeconómico}
%% ════════════════════════════════════════════════════════════════════════════

El presente trabajo analiza los determinantes de la inflación en Colombia durante
1975--2005 mediante un modelo econométrico de regresión múltiple en forma
\textbf{log-log}, estimado por Mínimos Cuadrados Ordinarios (MCO). Todos los
resultados han sido verificados directamente en Stata 17 bajo la guía de la
Profesora Rosalba Gil Galindo.

La muestra abarca tres regímenes macroeconómicos colombianos:

\begin{enumerate}[noitemsep]
  \item \textbf{Represión financiera / alta inflación} (1975--1989): inflación entre 16\%
        y 32\%, tasas de interés fijadas administrativamente.
  \item \textbf{Apertura y liberalización} (1990--1998): Constitución de 1991,
        independencia del Banco de la República, apertura al capital extranjero.
  \item \textbf{Metas de inflación} (1999--2005): inflación desciende de dos dígitos
        a un dígito; tasas de interés determinadas por el mercado.
\end{enumerate}

\begin{notabox}º
\begin{tabular}{ll}
\toprule
\textbf{Stata 19} & \textbf{R equivalente}\\
\midrule
\texttt{gen lninflacion = ln(inflacion)} & \texttt{datos\$lninflacion <- log(datos\$inflacion)}\\
\texttt{reg lninflacion lnpib lnti lnm} & \texttt{lm(lninflacion $\sim$ lnpib+lnti+lnm)}\\
\texttt{correlate lnpib lnti lnm} & \texttt{cor(datos[,c("lnpib","lnti","lnm")])}\\
\texttt{vif} & \texttt{car::vif(modelo)}\\
\texttt{reg lninflacion lnpib lnti} & \texttt{lm(lninflacion $\sim$ lnpib+lnti)}\\
\texttt{ovtest} & \texttt{lmtest::resettest(modelo,power=2:3)}\\
\bottomrule
\end{tabular}
\end{notabox}

%% ════════════════════════════════════════════════════════════════════════════
\section{Base de datos}
%% ════════════════════════════════════════════════════════════════════════════

\begin{longtable}{rrrrrrrr}
\caption{Colombia 1975--2005: datos originales y logarítmicos}\label{tab:datos}\\
\toprule
\textbf{Año}&\textbf{INF}&\textbf{PIB}&\textbf{TI}&\textbf{M}&
\textbf{ln\,INF}&\textbf{ln\,TI}&\textbf{ln\,M}\\
\midrule\endfirsthead
\multicolumn{8}{c}{\footnotesize\itshape(continuación)}\\\toprule
\textbf{Año}&\textbf{INF}&\textbf{PIB}&\textbf{TI}&\textbf{M}&
\textbf{ln\,INF}&\textbf{ln\,TI}&\textbf{ln\,M}\\\midrule\endhead
\bottomrule\endfoot
1975&17.70&405{,}108&23.80&51.9&2.8736&3.1697&3.9485\\
1976&25.70&424{,}263&22.40&82.7&3.2465&3.1091&4.4153\\
1977&28.40&441{,}906&22.90&103.5&3.3464&3.1311&4.6396\\
1978&18.80&479{,}335&25.90&134.9&2.9339&3.2542&4.9045\\
1979&28.80&505{,}119&36.50&167.6&3.3604&3.5973&5.1215\\
1980&26.80&525{,}765&41.50&214.3&3.2882&3.7257&5.3679\\
1981&26.40&537{,}736&35.70&259.7&3.2730&3.5757&5.5591\\
1982&24.00&542{,}836&37.70&325.7&3.1781&3.6290&5.7861\\
1983&16.60&551{,}380&37.30&408.9&2.8094&3.6181&6.0132\\
1984&18.30&569{,}855&34.60&503.8&2.9069&3.5441&6.2219\\
1985&22.50&587{,}561&35.90&642.2&3.1135&3.5797&6.4647\\
1986&21.00&621{,}781&32.10&788.5&3.0445&3.4696&6.6702\\
1987&24.00&655{,}164&34.20&1{,}048.3&3.1781&3.5326&6.9549\\
1988&28.10&681{,}791&31.80&1{,}318.5&3.3362&3.4592&7.1845\\
1989&26.10&705{,}088&34.00&1{,}702.1&3.2613&3.5264&7.4392\\
1990&32.40&24{,}030{,}173&38.20&2{,}141.0&3.4775&3.6428&7.6692\\
1991&26.80&31{,}130{,}592&37.10&2{,}842.0&3.2882&3.6133&7.9526\\
1992&25.10&39{,}730{,}752&27.20&4{,}013.1&3.2238&3.3033&8.2972\\
1993&22.60&52{,}271{,}688&25.60&5{,}214.8&3.1187&3.2426&8.5593\\
1994&22.60&67{,}532{,}862&29.50&6{,}419.0&3.1187&3.3844&8.7669\\
1995&19.50&84{,}439{,}102&32.00&7{,}717.8&2.9704&3.4657&8.9513\\
1996&21.60&100{,}711{,}389&31.43&8{,}992.8&3.0732&3.4475&9.1041\\
1997&17.70&121{,}707{,}501&24.09&10{,}948.0&2.8749&3.1812&9.3009\\
1998&16.70&140{,}483{,}322&32.55&10{,}526.5&2.8154&3.4826&9.2617\\
1999&9.23&151{,}565{,}005&21.60&12{,}856.8&2.2228&3.0741&9.4618\\
2000&8.75&174{,}896{,}258&12.14&16{,}720.8&2.1691&2.4966&9.7249\\
2001&7.65&188{,}558{,}786&12.47&18{,}737.0&2.0347&2.5232&9.8381\\
2002&6.99&203{,}451{,}414&8.92&21{,}635.6&1.9449&2.1884&9.9822\\
2003&6.49&228{,}516{,}603&7.80&24{,}918.3&1.8699&2.0541&10.1228\\
2004&5.50&239{,}439{,}697&7.80&29{,}113.7&1.7047&2.0541&10.2795\\
2005&4.85&251{,}722{,}954&7.03&26{,}954.4&1.5790&1.9509&10.2027\\
\end{longtable}

%% ════════════════════════════════════════════════════════════════════════════
\section{Punto 1: Especificación y Estimación del Modelo Log-Log}
%% ════════════════════════════════════════════════════════════════════════════

\subsection{Forma funcional del modelo}

El modelo multiplicativo en niveles es:
\begin{equation}
INF_t = \alpha\cdot PIB_t^{\beta_2}\cdot TI_t^{\beta_3}\cdot M_t^{\beta_4}\cdot e^{u_t}
\end{equation}

Tomando logaritmo natural a ambos lados se obtiene el modelo linealizado (\textbf{log-log}):

\begin{frmbox}{Ecuación Poblacional del Modelo Log-Log}
\begin{equation}
\ln(INF_t)=\underbrace{\beta_1}_{\text{intercepto}}
+\underbrace{\beta_2}_{\varepsilon_{INF,PIB}}\ln(PIB_t)
+\underbrace{\beta_3}_{\varepsilon_{INF,TI}}\ln(TI_t)
+\underbrace{\beta_4}_{\varepsilon_{INF,M}}\ln(M_t)
+u_t
\label{eq:modelo}
\end{equation}
$t=1975,\ldots,2005$;\quad$u_t\sim iid(0,\sigma^2)$;\quad$\beta_1=\ln\alpha$;\quad
cada $\beta_j$ es una \textbf{elasticidad parcial}.
\end{frmbox}

\subsection{Representación matricial}

Para $n=31$ observaciones, el modelo \eqref{eq:modelo} se escribe $\y=\X\bm\beta+\bm u$:

\[
\y_{31\times1}=\begin{bmatrix}\ln INF_{1975}\\\vdots\\\ln INF_{2005}\end{bmatrix},\qquad
\X_{31\times4}=\begin{bmatrix}
  1&\ln PIB_{1975}&\ln TI_{1975}&\ln M_{1975}\\
  \vdots&\vdots&\vdots&\vdots\\
  1&\ln PIB_{2005}&\ln TI_{2005}&\ln M_{2005}
\end{bmatrix},\qquad
\bm\beta=\begin{bmatrix}\beta_1\\\beta_2\\\beta_3\\\beta_4\end{bmatrix}
\]

\subsection{Estimación MCO}

Minimizando $\bm u'\bm u=(\y-\X\bm\beta)'(\y-\X\bm\beta)$ la condición de primer orden
produce el sistema de ecuaciones normales $\Xt\X\bm\beta=\Xt\y$, cuya solución única es:

\begin{frmbox}{Fórmula Fundamental MCO}
\begin{equation}
\bhat=(\Xt\X)^{-1}\Xt\y
\label{eq:mco}
\end{equation}
Existe si y sólo si $\mathrm{rango}(\X)=k=4$ (no hay multicolinealidad perfecta).\\
La varianza del estimador es $\Var(\bhat)=\sigma^2(\Xt\X)^{-1}$ (mínima por Gauss-Markov).
\end{frmbox}

\subsection{Resultados exactos de Stata 17}

\begin{statabox}
\begin{lstlisting}[style=Stata]
. reg lninflacion lnpib lnti lnm

      Source |       SS           df       MS      Number of obs =  31
-------------+----------------------------------   F(3, 27)      = 54.62
       Model |  7.67359572         3  2.55786524   Prob > F      = 0.0000
    Residual |  1.26440037        27  .046829643   R-squared     = 0.8585
-------------+----------------------------------   Adj R-squared = 0.8428
       Total |  8.93799609        30  .297933203   Root MSE      = .2164

 lninflacion | Coefficient  Std. err.      t    P>|t|   [95% conf. interval]
-------------+----------------------------------------------------------------
       lnpib |   .0391012   .0427228     0.92   0.368   -.0485587    .1267611
        lnti |   .8361166   .0949468     8.81   0.000    .6413019    1.030931
         lnm |  -.1026192   .0581286    -1.77   0.089   -.2218893    .0166509
       _cons |    .340877   .5566137     0.61   0.545   -.8011999    1.482954
\end{lstlisting}
\end{statabox}

\begin{table}[H]\centering
\caption{Resultados MCO --- Modelo Log-Log · Colombia 1975--2005}
\label{tab:mco}
\begin{tabular}{lrrrrcc}
\toprule
\textbf{Variable} & \textbf{Coeficiente} & \textbf{Std.\,Err.} & $t$ & $P>|t|$ &
IC 95\% inf & IC 95\% sup\\
\midrule
\rowcolor{gris_f}$\ln(PIB)$ & $\phantom{-}0.0391012$ & $0.0427228$ & $0.92$ & $0.368$
  & $-0.0486$ & $0.1268$\\
$\ln(TI)$ & $\phantom{-}0.8361166^{***}$ & $0.0949468$ & $8.81$ & $0.000$
  & $\phantom{-}0.6413$ & $1.0309$\\
\rowcolor{gris_f}$\ln(M)$ & $-0.1026192^{\dag}$ & $0.0581286$ & $-1.77$ & $0.089$
  & $-0.2219$ & $0.0167$\\
Constante & $\phantom{-}0.3408770$ & $0.5566137$ & $0.61$ & $0.545$
  & $-0.8012$ & $1.4830$\\
\midrule
\multicolumn{7}{l}{\small$n=31$\quad SCE$=7.6736$\quad SCR$=1.2644$\quad SCT$=8.9380$}\\
\multicolumn{7}{l}{\small$F(3,27)=54.62$, $p\approx0$\quad$R^2=0.8585$\quad
  $\bar R^2=0.8428$\quad Root\,MSE$=0.2164$}\\
\multicolumn{7}{l}{\footnotesize$^{***}p<0.001$;\quad$^{\dag}p<0.10$}\\
\bottomrule
\end{tabular}
\end{table}

\begin{resbox}
Los resultados del script R reproducen con precisión de 7 decimales la salida de Stata.
El modelo explica el \textbf{85.85\%} de la variación en $\ln(INF)$ y es globalmente
significativo ($F=54.62$, $p\approx0$). La variable de mayor peso individual es
$\ln(TI)$ con $t=8.81$ ($p\approx0$).
\end{resbox}

\subsection{Interpretación de las elasticidades}

En el modelo log-log: $\hat\beta_j=\dfrac{\partial\ln INF_t}{\partial\ln X_{jt}}
\approx\dfrac{\%\Delta INF}{\%\Delta X_j}$ (elasticidad parcial constante).

\begin{table}[H]\centering
\caption{Interpretación cuantitativa de las elasticidades estimadas}
\begin{tabularx}{\linewidth}{llX}
\toprule
\textbf{Parámetro} & \textbf{Valor} & \textbf{Lectura económica}\\
\midrule
$\hat\beta_2$ (lnpib) & $+0.0391$ &
  +1\% en el PIB $\Rightarrow$ +0.039\% en inflación, \textit{cet.\,par.}
  (efecto pequeño, no significativo). \\[4pt]
$\hat\beta_3$ (lnti) & $+0.8361^{***}$ &
  +1\% en la TI $\Rightarrow$ +0.836\% en inflación, \textit{cet.\,par.}
  (altamente significativo; refleja represión financiera en el período). \\[4pt]
$\hat\beta_4$ (lnm) & $-0.1026^{\dag}$ &
  +1\% en M $\Rightarrow$ $-$0.103\% en inflación, \textit{cet.\,par.}
  (signo atípico: efecto de colinealidad con lnpib; $p=0.089$ marginal). \\
\bottomrule
\end{tabularx}
\end{table}

\subsection{Términos clave del sistema financiero y política monetaria}

\begin{defbox}{Coeficiente de Efectivo ($ce$)}
$ce = C/D$, donde $C$ = efectivo en circulación y $D$ = depósitos bancarios.
Indica qué fracción de los medios de pago el público prefiere mantener en
efectivo. A mayor $ce$, menor multiplicador monetario:
$m = (1+ce)/(\rho+ce)$.
En Colombia, el alto $ce$ durante 1975--1989 restringió la expansión crediticia.
\end{defbox}

\begin{defbox}{Demasiado Grande para Quebrar (\textit{Too Big to Fail} --- TBTF)}
Institución financiera cuyo colapso generaría externalidades sistémicas tan
graves que el gobierno se ve políticamente obligado a rescatarla. Genera
\textbf{riesgo moral}: los directivos asumen riesgos excesivos al percibir que
las pérdidas serán socializadas. Regulado desde 2010 mediante recargos de
capital en Basilea III para bancos de importancia sistémica global (G-SIBs).
\end{defbox}

\begin{defbox}{Demasiado Grande para ser Rescatado (\textit{Too Big to Save})}
Situación donde el costo del rescate bancario excede la capacidad fiscal soberana.
Caso paradigmático: Irlanda (2010), cuyo rescate elevó la deuda pública del
$\approx$\,25\% al $\approx$\,120\% del PIB, forzando un rescate europeo.
\end{defbox}

\begin{defbox}{Demasiado Interconectado para Quebrar (\textit{Too Interconnected to Fail})}
El riesgo sistémico depende no sólo del tamaño sino de la \textbf{centralidad en
la red financiera}. Una institución mediana que actúa como nodo central de
compensación o contraparte dominante en mercados de derivados puede propagar
contagio a toda la red financiera ante su quiebra.
\end{defbox}

\begin{defbox}{Exceso de Reservas ($ER$)}
\[ER = R_{\text{total}} - R_{\text{requeridas}} = R_{\text{total}} - \rho\cdot D\]
Reservas mantenidas voluntariamente por encima del encaje mínimo. En recesiones
o períodos de incertidumbre, $ER$ sube porque los bancos prefieren liquidez a
crédito. El multiplicador monetario efectivo cae aunque la base $B$ sea elevada:
el banco central «empuja una cuerda» (\textit{pushing on a string}).
\end{defbox}

\begin{defbox}{Intervención en el Mercado Cambiario}
\textbf{Esterilizada:} el BC compra divisas (+$B$) y simultáneamente vende bonos
($-B$) para neutralizar el efecto monetario. Objetivo: controlar el tipo de cambio
sin alterar $M_s$. Efectividad limitada si hay movilidad perfecta de capitales.

\textbf{No esterilizada:} la compra de divisas se transmite íntegramente a la base
monetaria ($\Delta B > 0$). Más efectiva sobre el TC pero con consecuencias
inflacionarias (Colombia 1975--1989).
\end{defbox}

\begin{defbox}{Liberación Financiera (\textit{Financial Liberalization})}
Eliminación de controles sobre tasas de interés, crédito, flujos de capital y
acceso de banca extranjera. En Colombia, la apertura de 1990--1993 (Ley 45/1990,
Ley 35/1993) desmanteló la represión financiera previa. Aumenta la eficiencia
asignativa del sistema financiero pero también la vulnerabilidad a crisis externas
y ataques especulativos sobre la tasa de cambio.
\end{defbox}

\begin{defbox}{Multiplicador del Dinero ($m$)}
\begin{equation}
M_s = m\cdot B,\qquad m = \frac{1+ce}{\rho+ce},\qquad m>1\;\text{si}\;\rho<1
\end{equation}
$B$ = base monetaria; $\rho$ = tasa de encaje; $ce$ = coeficiente de efectivo.
Un banco central que expande $B$ en \$100 puede generar hasta \$100$\cdot m$ en
oferta monetaria amplia a través del proceso de creación de depósitos.
\end{defbox}

\begin{defbox}{Operaciones de Mercado Abierto (OMAs)}
Compraventa de valores públicos por el banco central en el mercado secundario.
Principal instrumento de política monetaria operativa:
\[\text{Compra bonos}\Rightarrow\uparrow B\Rightarrow\uparrow M_s\Rightarrow\downarrow i\Rightarrow\uparrow\pi
\qquad\text{Venta bonos}\Rightarrow\downarrow B\Rightarrow\downarrow M_s\Rightarrow\uparrow i\Rightarrow\downarrow\pi\]
El Banco de la República de Colombia las ejecuta mediante \textbf{repos} y
\textbf{repos en reversa} (SBRs) en subastas de expansión y contracción.
\end{defbox}

\begin{defbox}{Política de Encaje / Requisito de Reserva ($\rho$)}
\[\rho = \frac{R_{\text{req}}}{D}\in(0,1)\]
Fracción de los depósitos que los bancos comerciales deben mantener como reservas
(en bóveda o depositados en el banco central). En Colombia fue el principal
instrumento de política cuantitativa durante 1975--1990, antes de que se
desarrollaran los mercados de bonos internos. Variaciones en $\rho$ alteran
directamente el multiplicador $m$ sin modificar $B$.
\end{defbox}

\begin{defbox}{Prestamista de Última Instancia (\textit{Lender of Last Resort} -- LOLR)}
Función del banco central de proveer liquidez de emergencia a bancos
\textbf{solventes pero ilíquidos} para prevenir pánicos y corridas bancarias.
Regla de Bagehot (1873): \textit{prestar libremente, a tasa penalizadora, contra
colateral de buena calidad}. La ausencia de un LOLR creíble amplificó las crisis
bancarias latinoamericanas de los años 1990.
\end{defbox}

\begin{defbox}{Repo / Subasta de Expansión y Tasa Repo}
Operación en la que el BC compra temporalmente activos (títulos de deuda pública)
al sistema bancario comprometiéndose a revendérselos en fecha futura:
\[\text{Hoy: BC paga }P_0\qquad\text{Vencimiento: banco recompra por }P_0(1+i_\text{repo})\]
Inyecta liquidez (\textit{expansión}). La \textbf{Tasa Repo} es el instrumento
de política monetaria de referencia del Banco de la República desde 2001;
su variación transmite el estímulo o la contracción a toda la economía.
\end{defbox}

\begin{defbox}{Repo en Reversa / Subasta de Contracción y Tasa de Repos en Reversa}
Operación opuesta al repo: el BC vende temporalmente activos al sistema bancario,
retirando liquidez (\textit{contracción}, $\downarrow B$). La
\textbf{Tasa de Repos en Reversa} es el piso de la banda de tasas de interés
(rendimiento que el BC paga al sistema bancario por captar sus recursos
excedentarios \textit{overnight}).
\end{defbox}

\begin{defbox}{Represión Financiera (\textit{Financial Repression})}
Conjunto de regulaciones que mantienen las tasas de interés reales artificialmente
bajas (incluso negativas): topes administrativos a las tasas, encajes elevados,
compras forzosas de bonos públicos y restricciones a flujos de capital. Transfiere
riqueza implícitamente del sector privado al Estado (impuesto inflacionario).
Colombia operó bajo represión financiera durante 1975--1989, lo cual explica el
signo positivo de $\hat\beta_3$ en nuestro modelo.
\end{defbox}

%% ════════════════════════════════════════════════════════════════════════════
\section{Punto 2: Diagnóstico de Multicolinealidad Aproximada}
%% ════════════════════════════════════════════════════════════════════════════

\begin{defbox}{Multicolinealidad aproximada}
Existe cuando dos o más regresores mantienen una relación lineal aproximada (no
exacta). No viola ningún supuesto de Gauss-Markov pero \textbf{infla las
varianzas} de los estimadores: $\Var(\hat\beta_j)=\sigma^2 VIF_j/S_{jj}$.
Los estimadores MCO siguen siendo MELI pero con poca precisión individual.
La multicolinealidad es una característica de la muestra, no del modelo.
\end{defbox}

\subsection{2a. Verificación individual (t) y global (F)}

\subsubsection{Prueba global --- estadístico $F$}

\[H_0:\beta_2=\beta_3=\beta_4=0,\qquad
F=\frac{SCE/k_r}{SCR/(n-k)}=\frac{7.6736/3}{1.2644/27}=\frac{2.5579}{0.04683}=\mathbf{54.62}\]

$p$-valor $\approx 0.0000 < 0.05 \Rightarrow$ \textbf{RECHAZA} $H_0$: modelo globalmente significativo.

\subsubsection{Pruebas individuales --- estadístico $t$, $t_{0.025,27}=2.052$}

\begin{table}[H]\centering
\caption{Pruebas $t$ individuales y diagnóstico de multicolinealidad}
\begin{tabular}{lrrrrl}
\toprule
\textbf{Variable} & $\hat\beta$ & $SE$ & $t$ & $p$ & \textbf{Significancia}\\
\midrule
$\ln(PIB)$ & $+0.0391$ & $0.0427$ & $0.92$ & $0.368$ &
  \textcolor{rojo_a}{\textbf{NO significativa}} al 5\%\\
$\ln(TI)$  & $+0.8361$ & $0.0949$ & $8.81$ & $0.000$ &
  \textcolor{verde_o}{\textbf{Altamente significativa}} (***)\\
$\ln(M)$   & $-0.1026$ & $0.0581$ & $-1.77$ & $0.089$ &
  Marginal ($p<0.10$)\\
\midrule
\multicolumn{6}{l}{\small$R^2=0.8585$ (alto) + $F=54.62$ (significativo) + lnpib $p=0.368$ (no sig.)}\\
\multicolumn{6}{l}{\small$\Rightarrow$ \textbf{SÍNTOMA CLÁSICO de multicolinealidad}}\\
\bottomrule
\end{tabular}
\end{table}

\begin{alertbox}
\textbf{Síntoma clásico de multicolinealidad:} $R^2=0.8585$ muy alto, $F=54.62$
altamente significativo, pero $\ln(PIB)$ tiene $p=0.368$ (no significativo).
Esto ocurre porque lnpib y lnm están tan correlacionados ($r=0.9385$) que el
modelo no puede ``desenredar'' su contribución individual: infla sus errores
estándar y reduce sus estadísticos $t$ aunque su efecto conjunto es claro.
\end{alertbox}

\subsection{2b. Factor de Inflación de Varianza (VIF)}

\begin{frmbox}{Fórmulas VIF y Tolerancia}
\[VIF_j = \frac{1}{1-R_j^2} = \frac{1}{TOL_j},
\qquad\Var(\hat\beta_j)=\frac{\sigma^2}{S_{jj}}\cdot VIF_j\]
$R_j^2$ = $R^2$ de la regresión auxiliar de $X_j$ sobre los demás regresores.\\
\textbf{Reglas prácticas:} $VIF>10$: grave; $VIF\in(5,10)$: moderado-alto; $VIF<5$: aceptable.
\end{frmbox}

\begin{statabox}
\begin{lstlisting}[style=Stata]
. vif
    Variable |       VIF       1/VIF
-------------+----------------------
       lnpib |      8.72    0.114700
         lnm |      8.43    0.118636
        lnti |      1.61    0.620093
-------------+----------------------
    Mean VIF |      6.25
\end{lstlisting}
\end{statabox}

\begin{table}[H]\centering
\caption{VIF calculados manualmente y con Stata (coincidencia exacta)}
\begin{tabular}{lccccc}
\toprule
\textbf{Variable} & \textbf{Regresión auxiliar} & $R_j^2$ & $TOL_j$ & $VIF_j$ & \textbf{Diagnóstico}\\
\midrule
\rowcolor{naranja!15}$\ln(PIB)$ & lnpib$\sim$lnti+lnm & $0.8853$ & $0.1147$ & $8.72$ & Moderado-alto\\
$\ln(TI)$ & lnti$\sim$lnpib+lnm & $0.3799$ & $0.6201$ & $1.61$ & \textcolor{verde_o}{Bajo}\\
\rowcolor{naranja!15}$\ln(M)$ & lnm$\sim$lnpib+lnti & $0.8814$ & $0.1186$ & $8.43$ & Moderado-alto\\
\midrule
\multicolumn{5}{l}{Media VIF} & $6.25$\\
\bottomrule
\multicolumn{6}{l}{\footnotesize Umbral estricto: VIF$>10$. VIF$\approx8$--9 es ya una señal relevante de colinealidad.}
\end{tabular}
\end{table}

\begin{notabox}
Verificación manual: $VIF_{\ln PIB}=\dfrac{1}{1-0.8853}=\dfrac{1}{0.1147}=\mathbf{8.72}$
\quad$\checkmark$\quad Coincide exactamente con \texttt{estat vif} de Stata.
\end{notabox}

\subsection{2c. Test de Farrar-Glauber}

La matriz de correlación de los regresores (\texttt{correlate lnpib lnti lnm}):

\[
\mathbf{R}=\begin{pmatrix}
1.0000 & -0.6136 & 0.9385\\
-0.6136 & 1.0000 & -0.5959\\
0.9385 & -0.5959 & 1.0000
\end{pmatrix}
\]

\begin{statabox}
\begin{lstlisting}[style=Stata]
. correlate lnpib lnti lnm
             |    lnpib     lnti      lnm
-------------+---------------------------
       lnpib |   1.0000
        lnti |  -0.6136   1.0000
         lnm |   0.9385  -0.5959   1.0000
\end{lstlisting}
\end{statabox}

\begin{frmbox}{Estadístico $\chi^2$ de Farrar-Glauber}
\[H_0:\det(\mathbf{R})=1\quad\text{(no existe multicolinealidad)}\]
\[\chi^2_{FG}=-\!\left[n-1-\frac{2p+5}{6}\right]\!\ln\!\det(\mathbf{R})
\;\sim\;\chi^2\!\!\left(\frac{p(p-1)}{2}\right)\]
Con $p=3$ regresores y $n=31$: $gl=3(2)/2=3$.
\end{frmbox}

\[\det(\mathbf{R})=0.07396613\qquad\Rightarrow\qquad\ln\det(\mathbf{R})=-2.6041\]
\[\chi^2_{FG}=-\!\left[30-\frac{11}{6}\right](-2.6041)=-(28.1\overline{6})(-2.6041)=\mathbf{73.35}\]
\[p\text{-valor}\approx0.0000<0.05\;\Rightarrow\;\textbf{RECHAZA }H_0:\text{ multicolinealidad significativa}\]

\subsubsection{Correlaciones parciales (2ª etapa)}

\begin{table}[H]\centering
\caption{Correlaciones parciales y prueba $t$ (Farrar-Glauber)}
\begin{tabular}{lcccc}
\toprule
\textbf{Par} & $r_{ij\cdot\text{resto}}$ & $t$ & $p$ & \textbf{Conclusión}\\
\midrule
$r$(lnpib, lnti $\mid$ lnm) & $-0.1960$ & $-1.0385$ & $0.3082$ & No significativa\\
\rowcolor{rojo_a!10}$r$(lnpib, lnm $\mid$ lnti) & $\mathbf{+0.9033}$ & $\mathbf{+10.9436}$ &
  $\mathbf{0.0000}$ & \textbf{ALTAMENTE SIG.}\\
$r$(lnti, lnm $\mid$ lnpib) & $-0.0736$ & $-0.3835$ & $0.7044$ & No significativa\\
\bottomrule
\end{tabular}
\end{table}

\begin{resbox}
$r(\text{lnpib},\text{lnm}\mid\text{lnti})=0.9033$ con $t=10.94$, $p\approx0$.
PIB y Oferta Monetaria son casi linealmente dependientes aun controlando por TI:
ambas variables crecen en tendencia sincronizada durante todo el período.
\end{resbox}

\subsection{2d. Test de Theil}

\begin{frmbox}{Criterio de Theil}
$R_j^2$ (regresión auxiliar) $>$ $R^2$ (modelo principal) $\Rightarrow$ colinealidad problemática.
\end{frmbox}

\begin{table}[H]\centering
\caption{Test de Theil ($R^2_{\text{principal}}=0.8585$)}
\begin{tabular}{lccl}
\toprule
\textbf{Variable} & $R_j^2$ & $R_j^2>R^2_\text{princ.}$? & \textbf{Diagnóstico}\\
\midrule
\rowcolor{rojo_a!10}$\ln(PIB)$ & $0.8853$ & \textcolor{rojo_a}{\textbf{SÍ}} (0.8853$>$0.8585) & Colinealidad problemática\\
$\ln(TI)$ & $0.3799$ & No (0.3799$<$0.8585) & Sin problema\\
\rowcolor{rojo_a!10}$\ln(M)$ & $0.8814$ & \textcolor{rojo_a}{\textbf{SÍ}} (0.8814$>$0.8585) & Colinealidad problemática\\
\bottomrule
\end{tabular}
\end{table}

\begin{resbox}
\textbf{Conclusión unificada del Punto 2:} Los cuatro métodos convergen.
Existe multicolinealidad moderada-alta entre $\ln(PIB)$ y $\ln(M)$ ($r=0.9385$;
$VIF\approx8.7$; $\chi^2_{FG}=73.35$, $p\approx0$; Theil alerta en ambas).
La Tasa de Interés no presenta problemas de colinealidad ($VIF=1.61$).
\end{resbox}

%% ════════════════════════════════════════════════════════════════════════════
\section{Punto 3: Eliminación de Variables y Sesgo de Especificación}
%% ════════════════════════════════════════════════════════════════════════════

\subsection{Modelos reducidos (resultados exactos de Stata)}

\begin{statabox}
\begin{lstlisting}[style=Stata]
. reg lninflacion lnpib lnti          . reg lninflacion lnti lnm
  F(2,28)=74.72  R²=0.8422              F(2,28)=81.99  R²=0.8541
  lnpib: coef=-.0290308  t=-1.53        lnti: coef=.8190854  t=8.82***
  lnti:  coef=.8484537   t=8.64***      lnm:  coef=-.0545604  t=-2.19*
  VIF: lnpib=1.61  lnti=1.61            VIF: lnti=1.55  lnm=1.55

. reg lninflacion lnpib lnm
  F(2,28)=11.56  R²=0.4522
  lnpib: coef=-.0346351  t=-0.43
  lnm:   coef=-.1402957  t=-1.25
\end{lstlisting}
\end{statabox}

\begin{table}[H]\centering
\caption{Comparación: modelo completo vs.\ modelos reducidos}
\label{tab:reducidos}
\resizebox{\textwidth}{!}{%
\begin{tabular}{lccccc}
\toprule
& \textbf{Completo} & \textbf{Sin lnM} & \textbf{Sin lnPIB} & \textbf{Sin lnTI}\\
\midrule
$\ln(PIB)$ & $+0.0391\,(p\!=\!.37)$ &
  $\mathbf{-0.0290}\,[\textcolor{rojo_a}{\text{sesgo}}]$ & --- &
  $-0.0346\,[\textcolor{rojo_a}{\text{sesgo}}]$\\[2pt]
$\ln(TI)$ & $+0.8361^{***}$ & $+0.8485^{***}$ & $+0.8191^{***}$ & ---\\[2pt]
$\ln(M)$ & $-0.1026^{\dag}$ & --- &
  $-0.0546^{*}$ & $-0.1403\,[\textcolor{rojo_a}{\text{sesgo}}]$\\[4pt]
$R^2$ & $0.8585$ & $0.8422$ & $0.8541$ & $0.4522$\\
$\bar R^2$ & $0.8428$ & $0.8309$ & $0.8437$ & $0.4131$\\
Root\,MSE & $0.2164$ & $0.2244$ & $0.2158$ & $0.4182$\\
$F$ & $54.62$ & $74.72$ & $81.99$ & $11.56$\\
VIF\,máx & $8.72$ & $1.61$ & $1.55$ & $7.91$\\[4pt]
\midrule
MELI? & \YES & \NOO\,(sesgado) & \NOO\,(sesgado) & \NOO\,(sesgado)\\
\bottomrule
\end{tabular}}
\end{table}

\subsection{Fórmula del sesgo de especificación por omisión (OVB)}

\begin{frmbox}{Omitted Variable Bias (OVB)}
Si el modelo verdadero incluye $\ln M$ pero se estima sin él:
\begin{equation}
\E[\hat\alpha_2] = \beta_2 + \beta_4\cdot b_{42}
\label{eq:ovb}
\end{equation}
donde $b_{42}=\widehat{\Cov}(\ln PIB,\ln M)\,/\,\widehat{\Var}(\ln PIB)$.
El sesgo $\beta_4\cdot b_{42}\neq0$ requiere que la variable omitida sea
relevante ($\beta_4\neq0$) y correlacionada con las incluidas ($b_{42}\neq0$).
\end{frmbox}

\subsubsection{Cálculo numérico}

\[b_{42}=\frac{\widehat{\Cov}(\ln PIB,\ln M)}{\widehat{\Var}(\ln PIB)}
=\frac{6.3718}{9.3952}=0.6782\]

\[\text{Sesgo teórico}=\hat\beta_4\cdot b_{42}=(-0.1026192)\times0.6782=\mathbf{-0.06960}\]

\begin{table}[H]\centering
\caption{Verificación empírica de la fórmula OVB}
\begin{tabular}{lrl}
\toprule
Cantidad & Valor & Origen\\
\midrule
$\hat\beta_2$ en modelo pleno & $+0.0391012$ & MCO completo\\
$\hat\alpha_2$ en modelo sin lnM & $-0.0290308$ & MCO reducido\\
Diferencia observada & $-0.0681320$ & Diferencia\\
Sesgo teórico (fórmula \ref{eq:ovb}) & $-0.0696000$ & $(-0.1026)\times 0.6782$\\
\midrule
\multicolumn{3}{l}{\small Diferencia obs.\,($-0.0681$) $\approx$ sesgo teórico ($-0.0696$)
$\Rightarrow$ \textbf{validación empírica de OVB}}\\
\bottomrule
\end{tabular}
\end{table}

\begin{alertbox}
\textbf{Conclusión del Punto 3:} Eliminar $\ln(M)$ reduce el VIF de 8.72 a 1.61,
pero introduce un sesgo de $-0.070$ en el coeficiente del PIB que incluso cambia
su \textbf{signo} de positivo a negativo. Los estimadores del modelo reducido son
\textbf{sesgados e inconsistentes}. Bajo los supuestos clásicos, MCO en el modelo
completo sigue siendo MELI a pesar de la colinealidad.
\textbf{El remedio (sesgo) es peor que la enfermedad (colinealidad).}
\end{alertbox}

%% ════════════════════════════════════════════════════════════════════════════
\section{Punto 4: Coherencia Teórica del Modelo Estimado}
%% ════════════════════════════════════════════════════════════════════════════

\subsection{Marcos teóricos de referencia}

\textbf{Teoría Cuantitativa del Dinero (TCD):}
$MV=PQ\Rightarrow\ln P=\ln M+\ln V-\ln Q$.
Con $V$ estable: +1\% en $M\Rightarrow$ +1\% en $P$. Predicción: $\beta_4>0$.

\textbf{Modelo IS-LM:}
$\uparrow TI\Rightarrow$ contrae DA $\Rightarrow\downarrow P$.
Predicción: $\beta_3<0$.
$\uparrow PIB\Rightarrow$ expande DA $\Rightarrow\uparrow P$.
Predicción: $\beta_2>0$.

\textbf{Curva de Phillips:}
$\pi_t=\pi_t^e+\gamma(y_t-\bar y_t)+\varepsilon_t$.
Mayor output gap $\Rightarrow\uparrow\pi$. Predicción: $\beta_2>0$.

\subsection{Tabla de coherencia}

\begin{table}[H]\centering
\caption{Evaluación de coherencia teórica de los coeficientes estimados}
\label{tab:coh}
\begin{tabular}{lccccc}
\toprule
\textbf{Parámetro} & $\hat\beta$ & $p$ & \textbf{Signo esp.} & \textbf{Signo obs.} & \textbf{Coherente}\\
\midrule
$\hat\beta_2$ (lnpib) & $+0.0391$ & $0.368$ & $+$ (DA/Phillips) & $+$ & \YES\\
$\hat\beta_3$ (lnti)  & $+0.8361$ & $0.000$ & $-$ (IS-LM) & $+$ & \NOO\\
$\hat\beta_4$ (lnm)   & $-0.1026$ & $0.089$ & $+$ (TCD) & $-$ & \NOO\\
\bottomrule
\end{tabular}
\end{table}

\subsection{Análisis de las incoherencias}

\subsubsection{$\hat\beta_3=+0.8361$ --- lnTI positivo (esperado negativo por IS-LM)}

En Colombia durante 1975--1989, el sistema financiero operó bajo
\textbf{represión financiera}: las tasas de interés eran fijadas
administrativamente en niveles que, en términos reales, eran frecuentemente
negativos. Las tasas nominales subían \textit{porque} la inflación subía
(indexación reactiva), no como instrumento de control anti-inflacionario.
El coeficiente positivo refleja que la TI era un \textbf{indicador contemporáneo
de la inflación}, no un instrumento de contracción monetaria. La causalidad
econométrica captura la correlación $r(\ln INF,\ln TI)=0.9105$, no el canal de
política monetaria IS-LM. Sólo desde 1999 (metas de inflación con tasa repo
libre) la TI comenzó a actuar como instrumento. Una muestra dividida revelaría
$\hat\beta_3<0$ para el subperíodo 2000--2005.

\subsubsection{$\hat\beta_4=-0.1026$ --- lnM negativo (esperado positivo por TCD)}

Dos explicaciones complementarias:
\begin{enumerate}[noitemsep]
  \item \textbf{Colinealidad:} $r(\ln PIB,\ln M)=0.9385$ impide separar los
        efectos. La estimación de $\beta_4$ es altamente imprecisa ($p=0.089$).
  \item \textbf{Endogeneidad:} el Banco de la República expandía $M$
        precisamente cuando la inflación cedía (política expansiva en
        recesiones), generando correlación negativa espúrea entre $M$ y $\pi$.
\end{enumerate}

\begin{resbox}
\textbf{Conclusión del Punto 4:} Coherencia parcial. $\hat\beta_2>0$ es
consistente con DA/Phillips. $\hat\beta_3>0$ se explica por el contexto
institucional de represión financiera colombiana (1975--1989). $\hat\beta_4<0$
se atribuye principalmente a la multicolinealidad severa con lnpib y a posible
endogeneidad de la oferta monetaria. La variable de mayor solidez teórica y
estadística es $\ln(TI)$: $\hat\beta_3=0.8361$, $t=8.81$, $p\approx0$, domina
el poder explicativo ($R^2$ cae de 0.8585 a 0.4522 cuando se omite).
\end{resbox}

%% ════════════════════════════════════════════════════════════════════════════
\section{Supuestos MCO y diagnósticos adicionales}
%% ════════════════════════════════════════════════════════════════════════════

\begin{table}[H]\centering
\caption{Batería de pruebas sobre los supuestos MCO}
\begin{tabular}{llccc}
\toprule
\textbf{Supuesto} & \textbf{Test} & \textbf{Estadístico} & \textbf{$p$} & \textbf{Decisión}\\
\midrule
Normalidad & Shapiro-Wilk & $W=0.9730$ & $0.6036$ &
  \textcolor{verde_o}{Residuos normales \checkmark}\\
Homocedasticidad & Breusch-Pagan & (ver script) & --- &
  Ver script R\\
No autocorr. & Durbin-Watson & $DW=1.117$ & --- &
  \textcolor{rojo_a}{Posible autocorr. positiva \texttimes}\\
No autocorr. & Breusch-Godfrey(2) & (ver script) & --- &
  Ver script R\\
Especificación & RESET-Ramsey & (ver script) & --- &
  Ver script R\\
\bottomrule
\multicolumn{5}{l}{\small$DW=1.117 < d_L(31,3)\approx1.23$: señala autocorrelación positiva de 1er orden.}\\
\multicolumn{5}{l}{\small Frecuente en series de tiempo macroeconómicas. Corregir con errores HAC (Newey-West).}
\end{tabular}
\end{table}

%% ════════════════════════════════════════════════════════════════════════════
\section{Conclusiones generales}
%% ════════════════════════════════════════════════════════════════════════════

\begin{enumerate}
  \item \textbf{Estimación:} El modelo log-log ajusta bien
        ($R^2=0.8585$, $F=54.62$, $p\approx0$). Los coeficientes son elasticidades
        parciales verificadas exactamente en Stata. La variable dominante es
        $\ln(TI)$ con $\hat\beta_3=0.8361^{***}$.

  \item \textbf{Multicolinealidad:} Confirmada por los cuatro diagnósticos
        (VIF$\approx$8.7, síntoma clásico t/F, $\chi^2_{FG}=73.35$, Theil).
        La fuente es la alta correlación $r(\ln PIB,\ln M)=0.9385$. La TI
        no presenta problema.

  \item \textbf{Eliminación de variables:} Reduce el VIF de 8.72 a 1.61
        pero introduce un sesgo OVB de $-0.070$ que cambia el signo del
        coeficiente del PIB. MCO sigue siendo MELI en el modelo completo.
        No se recomienda eliminar variables.

  \item \textbf{Coherencia teórica:} Parcial. Los signos de lnTI y lnM son
        atípicos, explicados por represión financiera (lnTI) y colinealidad/
        endogeneidad (lnM). Se recomienda usar primeras diferencias,
        estimar por subperíodos o aplicar GMM con instrumentos para superar
        estas limitaciones.
\end{enumerate}

%% ════════════════════════════════════════════════════════════════════════════
\appendix
\section{Anexo A: Compendio de fórmulas econométricas}
%% ════════════════════════════════════════════════════════════════════════════

\begin{defbox}{Matrices del modelo MCO}
\begin{align*}
\text{Sombrero:}&\quad\mathbf{H}=\X(\Xt\X)^{-1}\Xt,\quad\mathbf{H}^2=\mathbf{H},\quad\text{tr}(\mathbf{H})=k\\
\text{Aniquilador:}&\quad\mathbf{M}=\mathbf{I}-\mathbf{H},\quad\hat{\bm u}=\mathbf{M}\y\\
\Var(\bhat)&=\sigma^2(\Xt\X)^{-1},\quad
\hat\sigma^2=\frac{\hat{\bm u}'\hat{\bm u}}{n-k}=\frac{SCR}{n-k}
\end{align*}
\end{defbox}

\begin{defbox}{Descomposición y métricas de ajuste}
\[SCT=SCE+SCR,\quad R^2=\frac{SCE}{SCT},\quad
\bar R^2=1-\frac{(1-R^2)(n-1)}{n-k}\]
\[F=\frac{SCE/k_r}{SCR/(n-k)}\sim F_{k_r,n-k},\quad
AIC=-2\ell+2k,\quad BIC=-2\ell+k\ln n\]
\end{defbox}

\begin{defbox}{Tabla resumen: instrumentos de política monetaria}
\begin{center}\small
\begin{tabular}{lllll}
\toprule
\textbf{Instrumento} & \textbf{Tipo} & $\Delta B$ & $\Delta i$ & $\Delta\pi$ esperado\\
\midrule
OMA compra bonos & Expansiva & $+$ & $\downarrow$ & $+$\\
OMA venta bonos  & Contractiva & $-$ & $\uparrow$ & $-$\\
Repo (SBR expansión) & Expansiva & $+$ & $\downarrow$ & $+$\\
Repo en reversa  & Contractiva & $-$ & $\uparrow$ & $-$\\
$\downarrow\rho$ (encaje) & Expansiva & $0\,(\uparrow m)$ & $\downarrow$ & $+$\\
$\uparrow\rho$ (encaje) & Contractiva & $0\,(\downarrow m)$ & $\uparrow$ & $-$\\
\bottomrule
\end{tabular}
\end{center}
\end{defbox}

%% ════════════════════════════════════════════════════════════════════════════
\section{Anexo B: Extracto del script R}
%% ════════════════════════════════════════════════════════════════════════════

\begin{lstlisting}[style=Rcode,caption={Fragmento clave del script R (réplica de Stata)}]
# gen lninflacion = ln(inflacion)
datos$lninflacion <- log(datos$inflacion)
datos$lnpib <- log(datos$pib)
datos$lnti  <- log(datos$ti)
datos$lnm   <- log(datos$m)

# Álgebra matricial MCO
y <- matrix(datos$lninflacion, ncol=1)
X <- cbind(1, datos$lnpib, datos$lnti, datos$lnm)
beta_hat <- solve(t(X)%*%X) %*% t(X)%*%y   # (X'X)^{-1} X'y

# reg lninflacion lnpib lnti lnm
modelo <- lm(lninflacion ~ lnpib + lnti + lnm, data=datos)
summary(modelo)   # Replica exacta de Stata

# correlate lnpib lnti lnm
cor(datos[,c("lnpib","lnti","lnm")])

# vif (Stata: estat vif)
car::vif(modelo)
# lnpib=8.72  lnm=8.43  lnti=1.61  Mean=6.25

# Farrar-Glauber
cor_reg <- cor(datos[,c("lnpib","lnti","lnm")])
detR    <- det(cor_reg)   # 0.07397
chi2_FG <- -(31-1-(2*3+5)/6)*log(detR)   # 73.35
pchisq(chi2_FG, df=3, lower.tail=FALSE)   # p ≈ 0

# Sesgo de omision OVB
b42   <- cov(datos$lnpib,datos$lnm)/var(datos$lnpib)  # 0.6782
sesgo <- coef(modelo)["lnm"]*b42                       # -0.0696

# reg lninflacion lnpib lnti  (sin lnM)
m_sinM <- lm(lninflacion ~ lnpib + lnti, data=datos)
car::vif(m_sinM)   # lnpib=1.61  lnti=1.61
\end{lstlisting}

%% ─────────────────────────────────────────────────────────────────────────────
\vspace{2cm}
\begin{center}
{\color{uptv}\rule{\linewidth}{2pt}}\vspace{2pt}
{\color{oro}\rule{\linewidth}{1.5pt}}
\vspace{4mm}

{\large\bfseries\color{uptv}Emanuel Quintana Silva}\\[4pt]
{\color{gris_t}Universidad Pedagógica y Tecnológica de Colombia (UPTC)}\\[3pt]
{\color{gris_t}Econometría}\\[6pt]
{\large\itshape\color{oro}Con agradecimiento a la Profesora Rosalba Gil Galindo}\\[4pt]
{\scriptsize\color{gris_t}Resultados verificados con Stata 17 $\cdot$ \today}

\vspace{4mm}
{\color{oro}\rule{\linewidth}{1.5pt}}\vspace{2pt}
{\color{uptv}\rule{\linewidth}{2pt}}
\end{center}

\end{document}
